% The poet travels to Rio de Janeiro for Carnival.
% He sits on the beach, observes the wildlife, and watches 
% a parade at the Sambodrome.

Or in January's city, \\
Plastic surgeon pecs, \\
Cosmetic sex \& preening.

--Coughing on the couch by day--
--The goblins eat their takeaway.--  

Booze brigade of fairybears march \\
Sharing, caring, underground. \\
Which echoed loudly like an anthem-- \\
Others in the tunnels joined them. \\
Contagious as a Flemish painting: \\
Pieter the Elder, maybe.

Cris provides a live translation \\
As the train arrives in station:

--I'd love to pour oil all over your hide-- \\
--Why? To watch your boomboom slide.--

\secdiv

Palm fronds gently dance a hula while you're flayed \& sizzled. \\
Towel-talking, loving language. \\
Wond'ring which were perfect popsongs of a poem, like \\
Frankie's ``Coke'', and which were perfect poems of popsongs? \\
Cohen's ``Suzy',' Joni's ``Ka{\textsci}.o{\textupsilon}t''-- \\
All of them (not ``Prufrock'') love-ode. 

While blue ships on a blue horizon \\
Read in cross-hatch rival currents. 

And shifting gaze past Sugarloaf, \\
Toward shaggy loafing lovers, \\
Sharing sodas, lemon máte, seasoned meat \\
It's mating season, meet'n'greet in pidgin Latin. \\
While blue ships on blue horizon \\
Drive toward trucks to trudge loads inland.

--Under pitched sunbrellas:-- \\
--Picturesque young Cinderellas.--

In portal promise of a paperback partake: \\
Omegaverse's wolfish sense of het-romantasy.

\secdiv

Forever \& always in the age of glass \\
The age of windowshopping \\
Love, \& sex, in crackt arcades \\
A tablet trance to faraway beach \\
In this, our Touring Twenties, \\
Testy, wanderyears sans wonder, \\
Picqued by pretty pictures, \\
Waxwing slain on false azure, \\
On subtlegridded portal promise, \\
Shallow light reflecting depth there.

\secdiv

That night catch bossa in the old cantina; \\
Cold beer lifting heat like evening \\
Breeze, \& kiddies craning cameras \\
Skin as brown as betel-nuts, and \\
Airy hymns from half-caught eras.

And midnight on the Sambodromo: \\
From our chairs we chatted \\
Charted floats slowprogress \\
Chanting crowd incanting chorus \\
Lowd \& rowdy \\
Cheering stars who flaunt \& shake \& \\
Flow to rhythmic current \\
Down the corso's \\
Shoreline, costumed cosplay torsos \\
Fest of fat and fleshy skin and \\
Winkless sleep slowsetting in, \\
The Sun King and his sinking Queen \\
Their crowns of rays \\
Their solar spokes arrayed \\
In radiating angle \\
Tended by a bend of angels \\
Featherwingèd, free of fetters \\
Wringing salt from festive tatters \\
Fire-singed and filing, singing \\
(Not trivial in tropic weathers) \\
Wending toward the barricades, the \\
Great brigades of cosplay tribes, the \\
Front-row seats secured by bribes.

"O céu vai clarear \\
Iluminar a zona oeste da cidade \\
E Deus vai desfilar \\
Pra ver o mago recriar a Mocidade 

A luz que nos chega da estrela primeira \\
Nascida do pó no Cruzeiro do Sul \\
Do plasma divino das mãos carpinteiras \\
Ressurge candeia no breu nesse azul \\
Será que o limbo da imaginação \\
Perverte a inteligência \\
O homem com sua ambição \\
Desconhece a razão desatina a Ciência \\
Será que há de ter carnaval, sem minha cadência? \\
Com alas em tom digital \\
No fim da existência \\
Me diz afinal quem há de arcar com as consequências?

Se a Mocidade sonhar \\
No infinito escrever \\
Versos a luz do luar, deixa! \\
Quando o futuro voltar \\
A juventude vai crer \\
Que toda estrela pode renascer

O verde adoecido da esperança \\
Ofega sobre o leito da cobiça \\
Quem vive pelo preço da cobrança \\
Derrama sua lágrima postiça \\
Fogo matando a floresta \\
Bicho morrendo no cio \\
Febre no pouco que resta \\
Secam as águas do rio \\
E a vida vai vivendo por um fio \\
Naveguei \\
No afã de me encontrar eu me emocionei \\
Lembrei da corda bamba que atravessei \\
São tantas as viradas desta vida \\
A mão que faz a bomba se arrepende \\
Faz o samba e aprende \\
A se entregar de corpo e alma na avenida."